% =====================================================================
% Section : Mécanisme de Popper en mode centralisé
% Généré à partir d'une lecture directe du code source de popper-core
% (symbolic/popper-core/popper/{loop,asp,generate,constrain,tester,util,core}.py
%  et symbolic/popper-core/popper/lp/alan.pl), intégré dans la plateforme
% FILP via engines/centralized/runner.py.
% =====================================================================

\section{Le mécanisme d'apprentissage centralisé de Popper}
\label{sec:popper-centralise}

Dans la configuration centralisée de la plateforme, l'apprentissage
n'implique aucune communication réseau ni agrégation : un unique
processus exécute directement la boucle d'apprentissage de
Popper~\cite{cropper2021popper} sur l'ensemble du jeu de données
(fichiers \texttt{bias.pl}, \texttt{bk.pl}, \texttt{exs.pl}), via
\texttt{popper.util.Settings} et \texttt{popper.loop.learn\_solution}.
Le moteur centralisé de la plateforme (\texttt{engines/centralized/}
\texttt{runner.py}) n'est qu'un mince processus enfant : il ne fait
qu'invoquer \texttt{learn\_solution(settings)} et sérialiser le résultat.
Ce choix d'un sous-processus n'est pas anodin : le mécanisme de
temporisation interne de Popper (voir~\S\ref{sec:popper-stop}) repose
sur \texttt{signal.alarm()}, qui n'est utilisable que depuis le thread
principal de l'interpréteur --- or les pages Streamlit de la plateforme
s'exécutent dans un thread secondaire. Isoler Popper dans son propre
processus lui redonne un thread principal légitime pour installer son
gestionnaire \texttt{SIGALRM}.

Popper implémente le paradigme \emph{generate, test, constrain}
(génère, teste, contraint) : à chaque itération, un programme candidat
est généré par résolution d'un problème de satisfaction de contraintes
encodé en \emph{Answer Set Programming} (ASP, via le solveur
Clingo~\cite{gebser2019multi}), puis testé contre les exemples positifs
et négatifs, et enfin --- si le résultat n'est ni une solution exacte ni
une impasse totale --- l'échec observé est traduit en une ou plusieurs
\emph{contraintes} qui sont réinjectées dans le solveur ASP pour
empêcher la regénération de programmes structurellement similaires.

\subsection{Vue d'ensemble : la boucle principale}

La fonction \texttt{popper()} (\texttt{loop.py}, lignes 195--252)
orchestre une double boucle : une boucle externe sur la \emph{taille}
du programme (nombre total de littéraux, tête et corps confondus) et,
pour chaque taille, une boucle interne qui énumère tous les programmes
ASP de cette taille jusqu'à épuisement (UNSAT), avant de passer à la
taille suivante. Les contraintes accumulées au fil de l'exécution ne
sont \textbf{jamais réinitialisées} entre deux tailles : le solveur ASP
est un unique objet \texttt{clingo.Control} persistant pour toute la
durée de la recherche, ce qui permet à Popper de tirer parti, à la
taille $n+1$, de tout ce qu'il a appris en éliminant des hypothèses à la
taille $n$.

\begin{algorithm}[H]
\caption{Boucle principale de Popper (\texttt{popper.loop.popper})}
\label{alg:popper-loop}
\begin{algorithmic}[1]
\Require Réglages $\mathit{settings}$ (biais, exemples, connaissances de fond, $\mathit{max\_literals}$)
\Ensure Un programme $P^\star$ (solution exacte ou meilleure approximation)
\State $\mathit{solver} \gets \textsc{ClingoSolver}(\mathit{settings})$ \Comment{charge \texttt{alan.pl} + \texttt{bias.pl}}
\State $\mathit{tester} \gets \textsc{Tester}(\mathit{settings})$ \Comment{moteur Prolog (SWI-Prolog via \texttt{pyswip})}
\State $\mathit{grounder} \gets \textsc{ClingoGrounder}()$
\State $\mathit{constrainer} \gets \textsc{Constrain}()$
\State $\mathit{best\_score} \gets \bot$
\For{$\mathit{size} = 1 \to \mathit{max\_literals}$} \Comment{itération sur la taille du programme}
    \State $\mathit{solver}.\textsc{UpdateNumberOfLiterals}(\mathit{size})$
    \While{\Call{True}{}}
        \State \Comment{--- GÉNÉRATION ---}
        \State $\mathit{model} \gets \mathit{solver}.\textsc{GetModel}()$
        \If{$\mathit{model} = \varnothing$}
            \State \textbf{break} \Comment{espace de recherche épuisé pour cette taille}
        \EndIf
        \State $(P, \mathit{before}, \mathit{min\_clause}) \gets \textsc{GenerateProgram}(\mathit{model})$
        \State \Comment{--- TEST ---}
        \State $(\mathit{tp}, \mathit{fn}, \mathit{tn}, \mathit{fp}) \gets \mathit{tester}.\textsc{Test}(P)$
        \State $\mathit{outcome} \gets \textsc{DecideOutcome}(\mathit{tp}, \mathit{fn}, \mathit{tn}, \mathit{fp})$
        \State $\mathit{score} \gets \mathit{tp} + \mathit{tn}$
        \If{$\mathit{best\_score} = \bot \lor \mathit{score} > \mathit{best\_score}$}
            \State $\mathit{best\_score} \gets \mathit{score}$
            \If{$\mathit{outcome} = (\textsc{All}, \textsc{None})$} \Comment{complet ET cohérent}
                \State \Return $P$ \Comment{solution exacte}
            \EndIf
            \State \textsc{Stats}.\textsc{RegisterBestProgram}($P$)
        \EndIf
        \State \Comment{--- CONTRAINTE ---}
        \State $R \gets \textsc{BuildRules}(\mathit{constrainer}, \mathit{tester}, P, \mathit{before}, \mathit{min\_clause}, \mathit{outcome})$
        \State $R_{\mathit{ground}} \gets \textsc{GroundRules}(\mathit{grounder}, R)$
        \State $\mathit{solver}.\textsc{AddGroundClauses}(R_{\mathit{ground}})$
    \EndWhile
\EndFor
\State \Return meilleur programme mémorisé (ou $\varnothing$ si aucun candidat n'a jamais été testé)
\end{algorithmic}
\end{algorithm}

Trois conditions d'arrêt coexistent (\S\ref{sec:popper-stop}) : la
découverte d'une solution exacte, l'épuisement de l'espace de recherche
jusqu'à $\mathit{max\_literals}$, ou l'expiration du budget de temps
alloué (\texttt{settings.timeout}, imposé par un \texttt{signal.alarm}
enveloppant tout l'appel à \texttt{popper()}).

\subsection{Phase de génération}
\label{sec:popper-generate}

\subsubsection{Encodage de l'espace des hypothèses en ASP}

L'espace des programmes candidats est encodé par un méta-programme ASP
fixe, \texttt{alan.pl} (chargé une fois par \texttt{ClingoSolver} au
démarrage, avec le fichier de biais spécifique au jeu de données). Les
principaux mécanismes qu'il définit sont~:

\begin{itemize}
  \item \textbf{Énumération des clauses.} Des règles de choix ASP
    (notation \texttt{\{...\}}) devinent, pour chaque clause $C \in
    \{0, \dots, \mathit{max\_clauses}-1\}$, quel littéral de tête
    (\texttt{head\_literal/4}) et quels littéraux de corps
    (\texttt{body\_literal/4}) elle contient, parmi les prédicats
    déclarés dans \texttt{bias.pl} (\texttt{head\_pred/2},
    \texttt{body\_pred/2}).
  \item \textbf{Taille du corps.} Le prédicat \texttt{body\_size(C,N)}
    compte les littéraux de corps distincts de la clause $C$ et une
    contrainte impose $1 \le N \le \mathit{max\_body}$.
  \item \textbf{Un seul littéral de tête par clause}, la clause $0$
    doit toujours exister, les clauses sont utilisées dans un ordre
    croissant, et les variables d'une clause sont utilisées dans un
    ordre croissant --- ces contraintes structurelles éliminent la
    redondance combinatoire due aux permutations équivalentes de
    clauses/variables.
  \item \textbf{Variables et arguments.} Les variables sont indexées
    $0, \dots, \mathit{max\_vars}-1$ ; un script Python embarqué
    (\texttt{\#script(python)}) énumère toutes les permutations de
    variables de l'arité requise pour chaque littéral.
  \item \textbf{Connexité, non-singularité, forme datalog.} Des
    contraintes éliminent les variables déconnectées, les variables
    singleton (sauf autorisation explicite) et les clauses non
    \emph{range-restricted}.
  \item \textbf{Typage et modes.} Les déclarations \texttt{type/2} et
    \texttt{direction/2} du fichier de biais sont traduites en
    contraintes \texttt{var\_type/3} et \texttt{direction\_/3},
    garantissant que deux occurrences de la même position de variable
    partagent un type cohérent, et qu'une notion de \emph{sûreté}
    modale (variable liée par une entrée \texttt{+} ou produite par un
    littéral entièrement en sortie) est respectée.
  \item \textbf{Ordre partiel \texttt{before/2} et \texttt{min\_clause/2}.}
    Précalculés pour encoder les contraintes d'ordre entre clauses
    (utile notamment pour la récursion et l'invention de prédicats), et
    réutilisés tels quels par le module de contrainte
    (\S\ref{sec:popper-constrain}).
\end{itemize}

\subsubsection{Contrôle incrémental de la taille du programme}

Plutôt que de re-sol\-veur/re-\allowbreak grounder tout l'espace des
hypothèses à chaque taille, Popper utilise un \emph{atome externe}
Clingo, ce qui permet de ne reconfigurer qu'une contrainte minuscule
entre deux itérations, tout en conservant l'ensemble des règles
grounded et des contraintes déjà apprises :
\[
  \texttt{\#external size\_in\_literals(n).}
\]
\[
  \texttt{:- size\_in\_literals(n), \#sum\{K+1,Clause : body\_size(Clause,K)\} != n.}
\]
À chaque changement de taille, \texttt{update\_number\_of\_literals(size)}
(1) désactive l'atome externe précédent, (2) instancie ce nouveau
fragment de programme ASP paramétré par la nouvelle taille via
\texttt{solver.ground(...)}, puis (3) active le nouvel atome externe
\texttt{size\_in\_literals(size)}. Le coût de cette opération est
indépendant de la taille de l'espace de recherche déjà exploré.

\subsubsection{De l'ensemble-réponse au programme Popper}

\texttt{ClingoSolver.get\_model()} demande au solveur incrémental le
prochain \emph{answer set} et retourne les symboles \og montrés \fg{}
(\texttt{head\_literal}, \texttt{body\_literal}, \texttt{before},
\texttt{min\_clause}, \texttt{direction\_}). La fonction
\texttt{generate\_program(model)} (\texttt{generate.py}) reconstruit
alors, pour chaque identifiant de clause présent dans le modèle, un
littéral de tête et un ensemble (frozenset) de littéraux de corps
Popper, avec leurs modes (\texttt{+}/\texttt{-}) dérivés des atomes
\texttt{direction\_}, ainsi que les structures auxiliaires
$\mathit{before}$ et $\mathit{min\_clause}$ qui seront transmises telles
quelles au module de contrainte.

\subsection{Phase de test}
\label{sec:popper-test}

Le test d'un programme candidat $P$ est délégué à un moteur Prolog réel
(SWI-Prolog, piloté depuis Python via \texttt{pyswip}). Au démarrage, le
\texttt{Tester} consulte dans l'ordre le fichier d'exemples, le fichier
de connaissances de fond, puis un méta-programme Prolog interne
(\texttt{test.pl}) qui définit notamment \texttt{success\_set/1}. Les
index des exemples positifs et négatifs sont extraits une fois pour
toutes (\texttt{self.pos}, \texttt{self.neg}).

\begin{itemize}
  \item Si toutes les clauses de $P$ sont \emph{séparables} (non
    récursives et n'impliquant aucun prédicat inventé), la couverture
    est calculée clause par clause puis unie --- une optimisation
    valide car des clauses séparables ne peuvent interagir entre elles.
  \item Sinon, le programme entier est asserté temporairement dans la
    base Prolog vivante (via un gestionnaire de contexte qui retire les
    clauses après usage), et \texttt{success\_set(Xs)} est interrogé
    une seule fois sur l'ensemble du programme.
\end{itemize}

La matrice de confusion $(\mathit{tp}, \mathit{fn}, \mathit{tn},
\mathit{fp})$ est alors calculée en comparant l'ensemble des exemples
couverts (prouvés par $P$) à \texttt{self.pos}/\texttt{self.neg} :

\[
  \mathit{tp} = |\mathit{pos} \cap \mathit{covered}|, \quad
  \mathit{fn} = |\mathit{pos} \setminus \mathit{covered}|, \quad
  \mathit{fp} = |\mathit{neg} \cap \mathit{covered}|, \quad
  \mathit{tn} = |\mathit{neg} \setminus \mathit{covered}|
\]

Deux prédicats auxiliaires évaluent une \emph{clause isolée} (et non le
programme entier) --- utilisés par la phase de contrainte pour un
raffinement plus fin quand $|P| > 1$ :
\[
  \textsc{IsTotallyIncomplete}(c) \iff \text{$c$ est séparable} \land
    \nexists\, x \in \mathit{pos} : x \in \mathit{success\_set}(\{c\})
\]
\[
  \textsc{IsInconsistent}(c) \iff \text{$c$ est séparable} \land
    \exists\, x \in \mathit{neg} : x \in \mathit{success\_set}(\{c\})
\]

\subsection{Décision de l'issue et score}
\label{sec:popper-outcome}

À partir de la matrice de confusion, \texttt{decide\_outcome} classe
séparément la couverture positive et la cohérence négative en trois
niveaux qualitatifs $\{\textsc{All}, \textsc{Some}, \textsc{None}\}$ :

\[
  \mathit{outcome}^+ =
  \begin{cases}
    \textsc{All}  & \text{si } \mathit{fn} = 0 \quad\text{(complet)} \\
    \textsc{None} & \text{si } \mathit{tp} = 0 \land \mathit{fn} > 0 \quad\text{(totalement incomplet)} \\
    \textsc{Some} & \text{sinon} \quad\text{(partiellement complet)}
  \end{cases}
\]
\[
  \mathit{outcome}^- =
  \begin{cases}
    \textsc{None} & \text{si } \mathit{fp} = 0 \quad\text{(cohérent)} \\
    \textsc{Some} & \text{sinon} \quad\text{(incohérent)}
  \end{cases}
\]

Le score utilisé pour retenir la \og meilleure hypothèse vue jusqu'ici \fg{}
est une simple exactitude non normalisée :
\[
  \mathit{score}(P) = \mathit{tp} + \mathit{tn}
\]
Une solution exacte est caractérisée par
$(\mathit{outcome}^+, \mathit{outcome}^-) = (\textsc{All}, \textsc{None})$,
c'est-à-dire $\mathit{fn} = \mathit{fp} = 0$ : le programme couvre tous
les exemples positifs et aucun exemple négatif. C'est la seule condition
qui provoque un retour immédiat de \texttt{popper()} --- toute autre
issue, même avec un score élevé, ne fait que mettre à jour le meilleur
programme mémorisé et alimente la phase de contrainte.

\subsection{Phase de contrainte}
\label{sec:popper-constrain}

C'est le cœur du mécanisme \og apprendre de ses échecs \fg{} de Popper :
chaque issue $(\mathit{outcome}^+, \mathit{outcome}^-)$ est associée à
un sous-ensemble de types de contraintes à ajouter à l'espace de
recherche ASP (table \texttt{OUTCOME\_TO\_CONSTRAINTS}, \texttt{loop.py}) :

\begin{table}[H]
\centering
\caption{Correspondance issue $\to$ contraintes générées}
\label{tab:outcome-constraints}
\begin{tabular}{@{}llll@{}}
\toprule
$\mathit{outcome}^+$ & $\mathit{outcome}^-$ & Interprétation & Contrainte(s) ajoutée(s) \\
\midrule
\textsc{All}  & \textsc{None} & solution exacte & (arrêt --- \textsc{Banish} en secours) \\
\textsc{All}  & \textsc{Some} & complet, incohérent & \textsc{Généralisation} \\
\textsc{Some} & \textsc{None} & incomplet, cohérent & \textsc{Spécialisation} \\
\textsc{Some} & \textsc{Some} & incomplet, incohérent & \textsc{Spécialisation} + \textsc{Généralisation} \\
\textsc{None} & \textsc{None} & totalement incomplet, cohérent & \textsc{Spécialisation} + \textsc{Redondance} \\
\textsc{None} & \textsc{Some} & totalement incomplet, incohérent & \textsc{Spécialisation} + \textsc{Redondance} + \textsc{Généralisation} \\
\bottomrule
\end{tabular}
\end{table}

La logique sous-jacente : une hypothèse \emph{déjà complète} mais
incohérente ne peut être corrigée qu'en la rendant plus spécifique ---
or Popper choisit ici de bannir ses \emph{généralisations} futures (une
généralisation resterait au moins aussi incohérente). Symétriquement,
une hypothèse \emph{incomplète} ne peut gagner en couverture qu'en
restant aussi générale ou en le devenant davantage --- ses
\emph{spécialisations} sont donc bannies (une spécialisation ne peut
que perdre de la couverture). Une hypothèse \emph{totalement}
incomplète (ne couvre aucun positif) est de plus un candidat naturel à
l'élimination \emph{redondante} : la garder ne peut qu'ajouter du bruit
au programme final.

Chaque générateur de contrainte (\texttt{Constrain.banish\_constraint},
\texttt{generalisation\_constraint}, \texttt{specialisation\_constraint},
\texttt{redundancy\_constraint}) construit, non pas des faits concrets,
mais des \emph{règles ASP abstraites} paramétrées par des
métavariables (\texttt{ConstVar}) représentant des numéros de clause ou
des identifiants de variable encore non fixés. Ces règles définissent un
prédicat auxiliaire \texttt{included\_clause/2} liant un \og gabarit \fg{}
de clause concret à tout numéro de clause qui l'instancierait, plus une
contrainte d'intégrité interdisant la re-génération de cette forme (ou
d'un sur-ensemble/sous-ensemble selon le type de contrainte). En
particulier :

\begin{itemize}
  \item \textsc{Banish} fixe le nombre exact de clauses du programme
    (via un littéral négatif sur \texttt{clause/1}) --- élimine
    précisément ce programme.
  \item \textsc{Généralisation} omet cette contrainte de cardinalité :
    la règle s'applique donc aussi à tout programme qui
    \emph{contiendrait} ce même jeu de clauses parmi d'autres --- ce
    qui, combiné à l'ordre imposé par le solveur, revient à interdire
    toute généralisation ultérieure de ce jeu de clauses.
  \item \textsc{Spécialisation} refixe la cardinalité exacte comme
    \textsc{Banish}, empêchant la re-dérivation de cette forme à la
    taille de programme courante --- les spécialisations, qui
    nécessitent des littéraux supplémentaires (donc une taille de
    programme strictement supérieure), sont ainsi naturellement
    repoussées, puis exclues, par la croissance de la contrainte
    \texttt{included\_clause}.
  \item \textsc{Redondance} calcule, pour chaque prédicat de tête
    apparaissant dans $P$, sa multiplicité (nombre de clauses) et sa
    dépendance récursive (via une clôture par point fixe), afin
    d'exclure de l'élagage tout prédicat impliqué dans une boucle de
    récursion, puis interdit toute forme conservant ce clause
    manifestement inutile.
\end{itemize}

Au-delà de la contrainte globale associée à l'issue du programme entier,
\texttt{build\_rules} affine encore le raisonnement \emph{clause par
clause} lorsque $|P| > 1$ : toute clause séparable jugée
\textsc{IsInconsistent} reçoit individuellement une contrainte de
\textsc{Généralisation} (même si l'issue globale du programme était
différente), et --- si toutes les clauses du programme sont séparables
--- toute clause jugée \textsc{IsTotallyIncomplete} reçoit
individuellement une contrainte de \textsc{Redondance}. Ce mécanisme
permet à Popper d'isoler et de pénaliser une clause fautive au sein d'un
programme à plusieurs clauses, sans nécessairement bannir tout le
programme.

\subsection{Phases de mise à terre (\emph{grounding}) et d'ajout}
\label{sec:popper-ground-add}

Les contraintes produites ci-dessus mélangent des littéraux ordinaires
et des \emph{méta-littéraux} (\texttt{AllDifferent}, $<$, $=$, $\ge$)
exprimés sur des métavariables symboliques. \texttt{ClingoGrounder}
résout ce petit problème de satisfaction de contraintes annexe : il
construit et résout un mini-programme ASP qui affecte à chaque
métavariable de clause une valeur dans $\{0,\dots,\mathit{max\_clauses}-1\}$
et à chaque métavariable de variable une valeur dans
$\{0,\dots,\mathit{max\_vars}-1\}$, sous les contraintes exprimées par
les méta-littéraux, et énumère \emph{toutes} les affectations valides.
Chaque affectation permet d'instancier (\emph{ground}) la règle
abstraite en une règle ASP concrète.

Enfin, \texttt{ClingoSolver.add\_ground\_clauses} utilise l'API
\emph{backend} bas niveau de Clingo pour injecter directement ces
règles concrètes dans le programme grounded du solveur, sans repasser
par un parseur ASP textuel --- chaque littéral est converti en atome
backend (mémoïsé pour éviter les doublons), puis chaque règle est
ajoutée comme contrainte d'intégrité ($\mathit{head} = \varnothing$) ou
comme règle définissant \texttt{included\_clause/2}. Ces règles
persistent dans l'unique objet \texttt{clingo.Control} pour le reste de
l'exécution, à toutes les tailles de programme suivantes.

\subsection{Conditions d'arrêt}
\label{sec:popper-stop}

Le mode centralisé n'ajoute \textbf{aucune} condition d'arrêt propre à
la plateforme : il expose directement la sémantique native de
\texttt{popper()}/\texttt{learn\_solution()} :

\begin{enumerate}
  \item \textbf{Solution exacte} --- retour immédiat dès qu'un
    programme atteint $(\textsc{All}, \textsc{None})$.
  \item \textbf{Espace de recherche épuisé} --- la boucle externe
    atteint $\mathit{max\_literals}$ (défaut~: 100) sans qu'aucune
    solution exacte n'ait été trouvée ; le meilleur programme mémorisé
    (par score $\mathit{tp}+\mathit{tn}$) est retourné.
  \item \textbf{Expiration du délai} --- \texttt{learn\_solution()}
    enveloppe l'appel à \texttt{popper()} dans un minuteur basé sur
    \texttt{signal.alarm()} (\texttt{settings.timeout}, défaut 600~s) ;
    à l'expiration, l'exécution est interrompue et le meilleur
    programme mémorisé jusque-là est retourné.
\end{enumerate}

Le nombre de \og rounds \fg{} n'a ici aucun sens fédéré : la
plateforme fixe simplement \texttt{number\_of\_rounds = 1} dans le
résultat sérialisé, puisqu'il n'existe qu'un seul processus
d'apprentissage, sans notion d'agrégation entre clients.

\subsection{Le fichier de biais~: un exemple concret}
\label{sec:popper-bias}

\begin{lstlisting}[language=Prolog, basicstyle=\small\ttfamily, breaklines=true, caption={Extrait d'un \texttt{bias.pl} (jeu de données \emph{trains})}]
max_clauses(4).
max_vars(3).
max_body(5).

head_pred(f,1).
body_pred(has_car,2).
body_pred(has_load,2).
body_pred(short,1).
body_pred(long,1).
body_pred(two_wheels,1).
body_pred(three_wheels,1).
body_pred(roof_open,1).
body_pred(roof_closed,1).
body_pred(zero_load,1).
body_pred(one_load,1).
body_pred(two_load,1).
body_pred(three_load,1).
body_pred(circle,1).
body_pred(triangle,1).
body_pred(rectangle,1).

type(f,(train,)).
type(has_car,(train,car)).
type(has_load,(car,load)).
type(short,(car,)).
% ... une ligne type/2 par predicat, donnant le type semantique
% de chaque argument.

:-
    clause(C),
    #count{V : var_type(C,V,train)} != 1.
\end{lstlisting}

Les trois premières lignes bornent directement la taille de l'espace de
recherche exploré par \texttt{ClingoSolver}/\texttt{ClingoGrounder}
(nombre maximal de clauses, de variables distinctes, et de littéraux de
corps par clause). \texttt{head\_pred/2} déclare l'unique prédicat
cible appris (ici \texttt{f/1}, classification d'un train), et
\texttt{body\_pred/2} énumère les 15 prédicats de fond utilisables dans
le corps des clauses. Les déclarations \texttt{type/2} imposent que
deux occurrences de la même position de variable partagent un type
sémantique cohérent (empêchant, par exemple, de confondre une variable
de type \texttt{train} avec une variable de type \texttt{car}). La
contrainte ASP finale, spécifique à ce jeu de données et \emph{non}
issue du méta-programme générique \texttt{alan.pl}, illustre que
\texttt{bias.pl} peut porter des règles d'intégrité arbitraires,
grounded et appliquées exactement comme celles d'\texttt{alan.pl} : ici,
elle force chaque clause apprise à lier exactement une variable de type
\texttt{train}, empêchant les clauses dégénérées ne portant sur aucun
train ou sur plusieurs simultanément.

% ---------------------------------------------------------------------
% Bibliographie minimale associée à cette section (à fusionner avec
% votre fichier .bib existant si les clés diffèrent).
% ---------------------------------------------------------------------
%
% @inproceedings{cropper2021popper,
%   title     = {Learning programs by learning from failures},
%   author    = {Cropper, Andrew and Morel, Rolf},
%   booktitle = {Machine Learning},
%   year      = {2021}
% }
%
% @article{gebser2019multi,
%   title   = {Multi-shot ASP solving with clingo},
%   author  = {Gebser, Martin and Kaminski, Roland and Kaufmann, Benjamin and Schaub, Torsten},
%   journal = {Theory and Practice of Logic Programming},
%   year    = {2019}
% }
